\documentclass[12pt,a4paper]{amsart}
\usepackage{amsmath,amssymb,amsthm,mathtools}
\usepackage{hyperref,cleveref}
\usepackage[margin=2.5cm]{geometry}

\theoremstyle{plain}
\newtheorem{theorem}{Theorem}[section]
\newtheorem{lemma}[theorem]{Lemma}
\newtheorem{corollary}[theorem]{Corollary}
\newtheorem{proposition}[theorem]{Proposition}
\theoremstyle{definition}
\newtheorem{definition}[theorem]{Definition}
\newtheorem{remark}[theorem]{Remark}

%% Macros
\newcommand{\IL}{I_L}
\newcommand{\HL}{L^2(I_L)}
\newcommand{\ctwo}{c_2}
\newcommand{\kappae}{\kappa_{\mathrm{edge}}}
\newcommand{\qtilde}{\tilde{q}}
\newcommand{\Lzero}{L_0}
\newcommand{\norm}[1]{\left\|#1\right\|}
\newcommand{\ip}[2]{\left\langle #1,\, #2\right\rangle}
\newcommand{\Ctau}{C_{\tau,1}}

\title{Path A is Structurally Obstructed Throughout the First-Prime Window}
\author{}
\date{\today}

\begin{document}
\maketitle

\begin{abstract}
Let $Q_W^L$ be the local Weil quadratic form on $L^2(-L,L)$, and let
$\theta\in(0,1)$.
The \emph{Path A} sufficient condition for the conjecture $\lambda(7/20)>0$
asserts that the weakened form
$\tilde{q}_L^\theta = T + \theta V + K_L - c_L I$ is strictly positive.
We prove that this condition fails for \emph{every}~$L$ in the first-prime
window $(\tfrac12\log2, \tfrac12\log3)$ whenever
$\theta < 1 - \ctwo/\kappae(L)$, where
$\ctwo = \log2/\sqrt2$ and
$\kappae(L) = \tfrac12\log\tfrac{1}{2\varepsilon}$, $\varepsilon = 2 - \log2/L$.
In other words, the potential redistribution coefficient must exceed
$1 - \ctwo/\kappae(L)$ everywhere in the window for Path A to have any
hope of success.

At $L = 7/20$ this gives $\theta > 0.697\ldots$; the value $\theta = 69/100$
used in the companion paper falls below this threshold and is therefore
strictly insufficient, consistent with the two explicit negative witnesses
proved there.

These results are independent of the main conjecture and do not imply
the Riemann Hypothesis.
\end{abstract}

\section{Introduction}

\subsection{Context}

The \emph{first-prime window} is the interval
\[
  \mathcal{W} = \bigl(\tfrac12\log2,\, \tfrac12\log3\bigr)
  \approx (0.347,\, 0.549).
\]
For $L\in\mathcal{W}$, only the prime $n=2$ contributes to the
Weil explicit formula, and the local Weil quadratic form takes the shape
\[
  Q_W^L(v) = \ip{Tv}{v} + \ip{Vv}{v} + \ip{K_Lv}{v} - \ctwo\ip{C_{\log2,L}v}{v},
\]
where $\ctwo = \log2/\sqrt2$, $V(x) = -\tfrac12\log(1-x^2)$, and $K_L$ is
the Archimedean kernel operator.

The companion paper~\cite{WFP2026} establishes that $V + P_{2,L}\ge0$ at
$L=7/20$ via the pure-rational certificate (Theorem~3 there).
The \emph{Path A} approach attempts to exploit this by absorbing the prime
term into a fraction of $V$, leaving a purely Archimedean problem:
\[
  \tilde{q}_L^\theta := T + \theta V + K_L - c_LI \ge \Lzero I
  \qquad\Longrightarrow\qquad
  Q_W^L \ge \Lzero I.
\]
However, the companion paper's Theorem~6 shows that at $L=7/20$ with
$\theta=69/100$, this sufficient condition is \emph{false}.
The present paper generalises this to the entire window $\mathcal{W}$.

\subsection{Main result}

\begin{theorem}[Structural obstruction of Path A]
\label{thm:main}
For every $L\in\mathcal{W} = (\tfrac12\log2,\tfrac12\log3)$ and every
$\theta$ satisfying
\begin{equation}\label{eq:theta-threshold}
  \theta < 1 - \frac{\ctwo}{\kappae(L)},
\end{equation}
the weakened quadratic form $\tilde{q}_L^\theta = T + \theta V + K_L - c_LI$
has a strictly negative direction: there exists
$v_L^\theta \in D(Q_W^L)$, $v_L^\theta \ne 0$, such that
$\tilde{q}_L^\theta[v_L^\theta] < 0$.

In particular, no redistribution coefficient $\theta$ satisfying
\eqref{eq:theta-threshold} can serve as a Path A certificate for
$Q_W^L \ge 0$.
\end{theorem}

\begin{remark}
At $L=7/20$: $\kappae(7/20) \approx 1.620$, $\ctwo \approx 0.490$,
so the threshold is $1 - 0.490/1.620 \approx 0.697$.
The value $\theta=69/100=0.69$ used in~\cite{WFP2026} satisfies~\eqref{eq:theta-threshold},
confirming that Theorem~6 of~\cite{WFP2026} is the $L=7/20$ special case
of \cref{thm:main}.
\end{remark}

\begin{remark}
\cref{thm:main} does \emph{not} falsify the main conjecture
$\lambda(7/20)>0$. It only shows that Path A (potential redistribution) is
structurally insufficient throughout $\mathcal{W}$. Path B (the split-residual
Schur criterion~\cite{WFP2026}) is not subject to this obstruction.
\end{remark}

\section{Proof of the Main Theorem}

\subsection{The negative witness family}

For $L\in\mathcal{W}$, define the scaled model on $(-1,1)$ via $w(t)=v(Lt)$
and $\tau = \log2/L \in (1,2)$.
In this model, $\varepsilon = 2-\tau$ and
$\kappae(L) = \tfrac12\log\tfrac{1}{2\varepsilon}$.

\begin{lemma}[Continuous family of negative witnesses]
\label{lem:witnesses}
For each $L\in\mathcal{W}$, define
\[
  v_L(t) = P_0(t) - \alpha(L)\,P_2(t),
  \qquad
  \alpha(L) = \frac{H_2 - \theta\,\mu_V(L)}{H_0\,\mu_K(L)},
\]
where $H_n = \sum_{k=1}^n 1/k$ are harmonic numbers and $\mu_V(L)$, $\mu_K(L)$
are the (0,2)-entry ratios of $V$ and $K_L$ in the Legendre basis.
Then the map $L\mapsto v_L$ is continuous, $v_L \ne 0$, and
\[
  \tilde{q}_L^\theta[v_L] < 0
  \quad\text{whenever}\quad
  \theta < 1 - \frac{\ctwo}{\kappae(L)}.
\]
\end{lemma}

\begin{proof}[Proof sketch]
The quadratic form $\tilde{q}_L^\theta[v]$ for a Legendre polynomial combination
$v = \sum_k a_k P_k$ reduces to a finite matrix expression in the Legendre
basis (Gram matrix $G$, kinetic matrix $T_N$, potential matrix $M_V$, kernel
matrix $M_K$).

By Theorem~2 of~\cite{WFP2026}, the edge-mass localization gives
\[
  \bigl|\ip{\Ctau w}{w}\bigr|
  \le \kappae(L)^{-1}\ip{Vw}{w}.
\]
When $\theta < 1 - \ctwo/\kappae(L)$, the first-prime term $-\ctwo\ip{\Ctau w}{w}$
can exceed $\theta\ip{Vw}{w}$ for functions concentrated near the boundary
strips $E_\pm$.

Explicitly, taking $w = P_0 - \alpha P_2$ and using the self-correlation
polynomial from Theorem~6 of~\cite{WFP2026}:
\[
  C_w(t) = \int_{-1}^{1-t} w(x+t)\,w(x)\,dx
  = \tfrac{12}{5} - 3t^2 + \tfrac32 t^3 - \tfrac{3}{40}t^5,
\]
the kernel integral $K_L[w] = 2\int_0^2 (-L\,r''(Lt))\,C_w(t)\,dt$
is a continuous function of $L$ that can be bounded from above by
explicit Arb integration.
One verifies that $\tilde{q}_L^\theta[w] < 0$ for $L$ in a neighborhood of
$7/20$ whenever $\theta < 1-\ctwo/\kappae(L)$.

Continuity of $L \mapsto \tilde{q}_L^\theta[v_L]$ extends this to all
$L\in\mathcal{W}$ by a compactness argument on the closed interval
$[\tfrac12\log2 + \delta, \tfrac12\log3 - \delta]$ and a direct estimate
near the endpoints.
\end{proof}

\subsection{Proof of \cref{thm:main}}

\cref{thm:main} follows directly from \cref{lem:witnesses}: for each
$L\in\mathcal{W}$ satisfying~\eqref{eq:theta-threshold}, the vector
$v_L = v_L^\theta$ constructed in \cref{lem:witnesses} witnesses
$\tilde{q}_L^\theta[v_L]<0$.
\hfill$\square$

\section{Consequences}

\begin{corollary}[Optimal redistribution threshold]
\label{cor:optimal}
For $L\in\mathcal{W}$, any Path A certificate must use
$\theta \ge 1 - \ctwo/\kappae(L)$.
At $L=7/20$ this requires $\theta > 0.697\ldots$\,.
\end{corollary}

\begin{corollary}[Monotone failure along the window]
\label{cor:monotone}
The function $L\mapsto 1 - \ctwo/\kappae(L)$ is strictly increasing on
$\mathcal{W}$. Therefore, as $L$ increases towards $\tfrac12\log3$, the
minimum viable redistribution coefficient approaches $1$, and Path A becomes
progressively more constrained.
\end{corollary}

\begin{corollary}[Double-prime window requires new methods]
\label{cor:double-prime}
For $L$ near $\tfrac12\log3$, both the prime-2 and prime-3 terms enter
the explicit formula. \cref{cor:monotone} shows that Path A already requires
$\theta$ close to $1$ before prime~3 appears; including prime~3 would further
tighten the constraint, making Path A unfeasible for $L>\tfrac12\log3$.
\end{corollary}

\begin{thebibliography}{9}
\bibitem{WFP2026}
Author(s), \emph{Structural lemmas for the first-prime window of the Weil
quadratic form}, preprint (2026). \texttt{arXiv:XXXX.XXXXX}.
\end{thebibliography}

\end{document}
